\documentclass[10pt,conference]{IEEEtran}
\IEEEoverridecommandlockouts
\usepackage{cite}
\usepackage{amsmath,amssymb,amsfonts}
\usepackage{graphicx}
\usepackage{textcomp}
\usepackage{xcolor}
\usepackage[hidelinks]{hyperref}
\begin{document}

\title{Hardware-Aware Quantum Bayesian Learner Compression:\\ Linking Human Belief Updating to NISQ Circuit Complexity\\
\large\normalfont\itshape IEEE Quantum Week (QCE26) --- Poster Proposal, Phase 2 (Sections 1--4)}

\author{
\IEEEauthorblockN{Mohammad Zoraiz}
\IEEEauthorblockA{
\textit{Pratt School of Engineering, Duke University} \\
Electrical Engineering, Physics, and Computer Science \\
Durham, NC, USA \\
mz248@duke.edu}
}

\maketitle

\begin{abstract}
This work connects quantum models of cognition with the engineering constraints of
near-term quantum hardware. Bayesian theories treat human learning as belief updating over
structured hypotheses, and quantum cognition represents belief states in Hilbert space with
evidence acting through quantum channels. We instantiate a small quantum Bayesian learner: a
belief state on a two-qubit register, a prior, and a stimulus-dependent evidence channel that
updates the belief on a synthetic similarity-structured categorization task inspired by the
simplicity principle of Pothos and Chater, with the posterior category read out by
measurement. To tie this cognitive picture to real devices, we make the learner a
hardware-efficient variational circuit and add a hardware-aware penalty equal to the number
of two-qubit gates that survive transpilation to IBM's FakeManila coupling map, then compress
the circuit by greedy structured pruning of individual Heisenberg interaction terms. The
central question is where the capacity boundary lies: how much entangling structure can be
removed before the learner can no longer represent the task. Across three difficulty levels
and five seeds we trace accuracy--complexity frontiers and find that structured compression is
essentially free, and on the hardest task beneficial, while removing all entanglers collapses
the learner to chance. Learned masks beat random masks at every matched budget. Finally, on a
physical IBM device the full circuit collapses to near chance while the compressed circuit
retains accuracy: on real NISQ hardware, hardware-aware compression is not merely economical
but necessary.
\end{abstract}

\begin{IEEEkeywords}
Quantum cognition, Bayesian learning, variational quantum circuits, hardware-efficient ansatz,
NISQ devices, circuit compression, quantum machine learning, quantum channels.
\end{IEEEkeywords}

\section{Poster Authors and Contact}
\textbf{Mohammad Zoraiz} (sole author; \emph{corresponding author} and \emph{presenter}).
Pratt School of Engineering, Duke University, Durham, NC, USA. Email:
\texttt{mz248@duke.edu}. ORCID: \emph{(optional --- add if available).} The title above is the
poster title (CFP item~1); this section is the author/contact information (CFP item~2); the
abstract above is the 200--250\,word poster abstract (CFP item~3); the relevance statement
below is CFP item~4.

\section{Poster Relevance to QCE26 Topics}
This poster is directly relevant to several QCE26 tracks. It is primarily a contribution to
\emph{Quantum Machine Learning} and \emph{Quantum Algorithms \& Information}: the learner is a
trainable, parameterized variational quantum circuit optimized by exact gradients, and the
compression study is a task-specific instance of quantum architecture search. It speaks to
\emph{Quantum Computing} and \emph{Quantum Systems Engineering} through its hardware-aware
cost---a real post-transpile two-qubit gate count on an IBM coupling map---and its validation
on a physical IBM Quantum device, directly addressing the dominance of two-qubit error on NISQ
hardware. The density-operator and quantum-channel formulation, with a maximally mixed prior
and a non-unital evidence update, connects to \emph{Quantum Generative AI} themes such as
quantum Boltzmann-style mixed-state models, while the per-term pruning of entangling
interactions contributes to \emph{Quantum Software Engineering} via circuit compression,
transpilation, and architecture optimization. By linking belief updating in quantum cognition
to circuit capacity and validating the resulting trade-off on real hardware, the work offers a
cross-disciplinary perspective of broad interest to the quantum computing and engineering
community. The preliminary poster follows on the next page (CFP Section~5); the accompanying
two-page extended abstract is submitted as a separate PDF (CFP Section~6).

\end{document}
